\section{A large clique}\label{sec:clique}

We continue with the graph from Section~\ref{sec:reduction}, and write
\(t=\omega(G)\) for its maximum clique size.  The next result supplies the
clique around which we will organize the rest of the proof.

\begin{theorem}\label{clq:bound}
Let \(k\ge3\) and \(2\le a\le k\) be integers, and let \(G\) be a finite
simple graph such that
\[
 |V(G)|=ka+1,\qquad \alpha(G)\le a,\qquad C_{k+1}\not\subseteq G.
\]
Suppose also that
\[
 |N_G[I]|\ge k|I|+1
 \quad\text{for every nonempty independent set }I\subseteq V(G).
\]
Then its maximum clique size \(t\) satisfies
\begin{equation}\label{clq:bound-equation}
 \max\{3,\lfloor k/2\rfloor\}\le t\le k.
\end{equation}
\end{theorem}

Lemma~\ref{red:expansion} gives the expansion hypothesis for our graph.
We first obtain a triangle.  For larger \(k\), we then find an induced
subgraph in a single distance layer with enough internal expansion to
force a long path.  Tree paths above that layer will close the path into
a cycle on exactly \(k+1\) vertices. This layer-and-tree strategy has its
antecedent in Erd\H{o}s, Faudree, Rousseau and Schelp
\cite[Section~3]{ErdosFaudreeRousseauSchelp1978}; the quantitative layer
and coloured-path statements used here are proved below.

\begin{lemma}\label{clq:triangle}
Under the hypotheses of Theorem~\ref{clq:bound}, we have \(3\le t\le k\).
\end{lemma}

\begin{proof}
A clique on \(k+1\) vertices contains a \(C_{k+1}\), so \(t\le k\).
For the other inequality, put \(r_b=R(K_3,K_b)\).  The elementary
Ramsey recurrence gives
\[
 r_2=3,\qquad r_b\le b+r_{b-1}\quad(b\ge3),
 \qquad r_b\le\frac{b(b+1)}2.
\]
Indeed, in a graph on \(b+r_{b-1}\) vertices, choose a vertex \(v\).
If \(v\) has at least \(b\) neighbors, either two of them are adjacent,
giving a triangle, or there is an independent set of size \(b\).
Otherwise \(v\) has at least \(r_{b-1}\) nonneighbors; among them there
is a triangle or an independent set of size \(b-1\), to which \(v\)
can be added.  Induction gives the displayed bound on \(r_b\).

Since \(a\le k\) and \(k\ge3\), we have \(a+3\le2k\), and hence
\[
 R(K_3,K_{a+1})\le\frac{(a+1)(a+2)}2\le ka+1=|V(G)|.
\]
The alternative of an independent set of size \(a+1\) is excluded,
so \(G\) contains a triangle.
\end{proof}

For \(k\le7\), Lemma~\ref{clq:triangle} already proves
Theorem~\ref{clq:bound}.  We therefore consider \(k\ge8\), and put
\(s=\lfloor k/2\rfloor\ge4\).  The following proposition isolates
the information needed from the distance layers.  It does not require
the whole graph to be connected.

\begin{proposition}\label{clq:layer-interface}
Let \(k\ge8\) be an integer, set \(s=\lfloor k/2\rfloor\), and let \(G\) be a
nonempty finite simple graph with \(\alpha(G)\le k\) and
\[
 |N_G[I]|\ge k|I|+1
 \quad\text{for every nonempty independent set }I\subseteq V(G).
\]
Suppose that \(G\) contains no \(K_s\).  Fix a vertex \(r\) and a
breadth-first spanning tree \(T\) of its component, rooted at \(r\).
Then there exist \(i\in\{1,\ldots,s\}\) and an induced subgraph \(H\)
of \(G\), all of whose vertices have depth \(i\) in \(T\), such that
\(H\) is connected and nonbipartite, and
\begin{equation}\label{clq:H-properties}
 \begin{gathered}
 |V(H)|\ge k,\qquad \delta(H)\ge s,\\
 |N_H[I]|\ge s|I|
 \quad(I\ne\varnothing\text{ independent in }H).
 \end{gathered}
\end{equation}
The last inequality is strict whenever \(k\) is odd or
\(H-N_H[I]\) is nonempty.
\end{proposition}

\begin{proof}
Let \(D_j\) be the set of vertices at distance \(j\) from \(r\).
These are exactly the depth layers of \(T\).  Layers beyond the
maximum depth are empty, and we use \(\alpha(\varnothing)=0\).
Every edge in the root component joins layers whose indices differ
by at most one.  Thus maximum independent sets chosen in layers of
one parity can be united, giving an independent set of size
\[
 p_i:=\sum_{\substack{0\le j\le i\\j\equiv i\pmod2}}
             \alpha(D_j).
\]
In particular \(p_i\le k\) and \(p_0=1\).  Applying expansion to a
singleton gives \(\delta(G)\ge k\), so \(D_1\ne\varnothing\).
Consequently every union defining \(p_i\) is nonempty: it contains
either the root or an independent vertex from \(D_1\).

For \(1\le i\le s\), the closed neighborhood of the independent
union defining \(p_{i-1}\) is contained in
\(D_0\cup\cdots\cup D_i\).  No vertex in another component is
adjacent to this union.  Expansion and \(|D_0|=1\) give
\begin{equation}\label{clq:layer-sum}
 \sum_{j=1}^i|D_j|\ge kp_{i-1}.
\end{equation}
Also, the two parity sums partition the indices from zero through
\(i\), so
\begin{equation}\label{clq:parity-sum}
 \sum_{j=1}^i\alpha(D_j)=p_i+p_{i-1}-1.
\end{equation}

Suppose first that \(k=2s\).  If every nonempty layer \(D_j\),
\(1\le j\le s\), satisfied \(|D_j|<s\alpha(D_j)\), then
\eqref{clq:layer-sum} and \eqref{clq:parity-sum} would imply
\[
 2sp_{i-1}\le\sum_{j=1}^i|D_j|
   <s(p_i+p_{i-1}-1)\qquad(1\le i\le s).
\]
The strict inequality holds even when some layers are empty, since
every sum contains the nonempty layer \(D_1\).
It follows that \(p_i>p_{i-1}+1\), and therefore, by integrality,
\(p_i\ge p_{i-1}+2\).  Iterating to \(i=s\) gives
\(p_s\ge2s+1>k\), a contradiction.  Hence some nonempty layer
\(D_i\), \(1\le i\le s\), satisfies
\begin{equation}\label{clq:even-layer}
 |D_i|\ge s\alpha(D_i).
\end{equation}

Now suppose that \(k=2s+1\).  If \(|D_j|\le s\alpha(D_j)\) for
every \(1\le j\le s\), including empty layers, the same two
identities give
\[
 (2s+1)p_{i-1}\le s(p_i+p_{i-1}-1),
 \qquad
 p_i\ge p_{i-1}+1+\left\lceil\frac{p_{i-1}}s\right\rceil.
\]
Starting from \(p_0=1\), each increment is at least two, so
\(p_{s-1}\ge2s-1>s\).  The last increment is consequently at least
three, giving \(p_s\ge2s+2>k\), again a contradiction.  Thus some
layer \(D_i\), \(1\le i\le s\), satisfies
\begin{equation}\label{clq:odd-layer}
 |D_i|>s\alpha(D_i).
\end{equation}
This layer is necessarily nonempty.

We have found a layer of sufficiently large order relative to its
independence number.  Choose within it a nonempty induced subgraph
\(H\) of minimum order satisfying the corresponding inequality
\eqref{clq:even-layer} or \eqref{clq:odd-layer}.  Since \(G\) has
no \(K_s\), such an \(H\) cannot be complete: a complete graph
satisfying either inequality would have at least \(s\) vertices.
Thus \(\alpha(H)\ge2\).  For even \(k\) this gives
\(|V(H)|\ge2s=k\); for odd \(k\) it gives
\(|V(H)|>2s\), hence \(|V(H)|\ge2s+1=k\).

The graph \(H\) is connected.  Otherwise each component \(H_j\)
would be a proper nonempty induced subgraph and would fail the
chosen inequality.  Since independence number is additive across
components, summing these failures would give
\[
 \begin{cases}
 |V(H)|<s\alpha(H),&k\text{ even},\\
 |V(H)|\le s\alpha(H),&k\text{ odd},
 \end{cases}
\]
contrary to its choice.  Nor is \(H\) bipartite: a bipartition
would give \(\alpha(H)\ge |V(H)|/2\), incompatible with
\(|V(H)|\ge s\alpha(H)\) and \(s\ge4\).

Let \(I\) be a nonempty independent set in \(H\), and put
\(R=H-N_H[I]\).  An independent set in \(R\) can be united with
\(I\), so
\[
 \alpha(R)\le\alpha(H)-|I|.
\]
If \(R\ne\varnothing\), minimality gives
\(|V(R)|<s\alpha(R)\) for even \(k\), and
\(|V(R)|\le s\alpha(R)\) for odd \(k\).  Comparing with the
defining inequality for \(H\), one comparison is strict in each
parity.  Therefore
\[
 |N_H[I]|=|V(H)|-|V(R)|
   >s\bigl(\alpha(H)-\alpha(R)\bigr)\ge s|I|.
\]
If \(R=\varnothing\), instead use
\(|N_H[I]|=|V(H)|\ge s\alpha(H)\ge s|I|\); the first
inequality is strict for odd \(k\).

Finally, take \(I=\{v\}\).  When \(H-N_H[v]\) is nonempty, the
strict neighborhood bound gives \(d_H(v)+1>s\), hence
\(d_H(v)\ge s\).  When it is empty, the order bound gives
\(d_H(v)=|V(H)|-1\ge k-1\ge s\).  This proves all the asserted
properties of \(H\).
\end{proof}

We now have a connected, nonbipartite graph \(H\) inside a single
depth layer \(D_i\), with \(1\le i\le s\), satisfying
\eqref{clq:H-properties}.  Under the assumption that \(G\) has no
\(K_s\), the same is true of \(H\).  The next argument uses these
properties to construct a path in \(H\) whose ends can be joined
through the tree above \(D_i\), producing the forbidden cycle.
